\documentclass[12pt]{article}
\usepackage{../sty/solspt}

\title{\sffamily\bfseries{NTD de Física 1º bimestre}}
\author{Samuel de Araújo Brandão}
\date{\today}

\begin{document}
\maketitle
Caro professor Isaías, estou encaminhando o NTD de física por meio deste documento escrito em \LaTeX,
pois considero que digito muito mais rápido que escrevo no caderno. Além de ficar mais legível e
facilitar o compartilhamento.

Quero esclarecer também que utilizei muita Inteligência Artifical. Tenho pouco tempo fora das aulas,
já que estudo muito para olimpíadas de matemática, já que em brever muitas delas ocorreram. Porém, pode
ter certeza que, quando eu for estudar melhor para a VP, buscarei compreender melhor as questões para
realizar as verificações parciais mais preparado.

\tableofcontents

\clearpage

\section{Questão 5}
\begin{enumerate}
	\item[(a)] Como os carros vão no mesmo sentido,
	      \[
		      v_{A/B}=v_A-v_B=80-90=\boxed{-10\ \text{km/h}}.
	      \]
	      (Isto significa que, para quem está em \(B\), \(A\) “anda para trás” a \(10\ \text{km/h}\).)

	\item[(b)] A velocidade de aproximação é
	      \[
		      v_{B/A}=v_B-v_A=90-80=10\ \text{km/h}.
	      \]
	      Com \(d=500\ \text{m}=0.5\ \text{km}\),
	      \[
		      t=\frac{d}{v_{B/A}}=\frac{0.5}{10}=0.05\ \text{h}=3\ \text{min}=\boxed{180\ \text{s}}.
	      \]
\end{enumerate}

\clearpage

\section{Questão 8}
Seja \(t=0\) no instante em que \(A\) passa pelo ponto e \(B\) parte.

\[
	x_A(t)=v_A t=10t
	\qquad\text{e}\qquad
	x_B(t)=\frac12 at^2=\frac12\cdot 4\,t^2=2t^2.
\]

No encontro: \(x_A=x_B\):
\[
	2t^2=10t \;\Rightarrow\; t(2t-10)=0 \;\Rightarrow\; t=5\,\text{s}.
\]

Distância percorrida por \(B\) até alcançar \(A\):
\[
	x_B(5)=2\cdot 5^2= \boxed{50\,\text{m}}.
\]

\clearpage

\section{Questão 12}
\begin{enumerate}
	\item[(a)] Queda livre a partir do repouso, de \(h=1{,}25\,\text{m}\) (adote \(g=10\,\text{m/s}^2\)):
	      \[
		      v=\sqrt{2gh}=\sqrt{2\cdot 10\cdot 1{,}25}=\sqrt{25}= \boxed{5\,\text{m/s}}.
	      \]

	\item[(b)] A torneira solta \(2\) gotas por minuto \(\Rightarrow\) intervalo de soltura
	      \[
		      \Delta t_0=\frac{60}{2}=30\,\text{s}.
	      \]
	      Como todas as gotas caem da mesma altura e do repouso, o tempo de queda é o mesmo para todas, então o intervalo entre impactos é
	      \[
		      \boxed{\Delta t=\Delta t_0=30\,\text{s}}.
	      \]
\end{enumerate}

\clearpage

\section{Questão 15}

\clearpage

\section{Questão 20}

\clearpage

\section{Questão 21}

\end{document}
